\chapter{Menelaus's Theorem}

\begin{definition}[Signed Length of a Segment]
    \label{def:Signed_Length}
    For any two distinct points A and B on an oriented line, the signed length $\overline{AB}$ is the distance between A and B, taken as positive if the direction from A to B is the same as the orientation of the line, and negative otherwise. For any three collinear points A, B, and C, this implies $\overline{AB} = -\overline{BA}$ and $\overline{AC} = \overline{AB} + \overline{BC}$. A ratio $\frac{\overline{AP}}{\overline{PB}}$ is positive if P is between A and B, and negative otherwise.
\end{definition}

\begin{definition}[Transversal of a Triangle]
    \label{def:Transversal}
    A transversal of a triangle $\triangle ABC$ is a line that intersects the lines containing the sides of the triangle, at points distinct from the vertices A, B, and C.
\end{definition}

\begin{lemma}[AA Similarity Criterion]
    \label{lem:AA_Similarity}
    If two angles of one triangle are congruent to two angles of another triangle, then the two triangles are similar. This implies that the ratio of corresponding side lengths is constant.
\end{lemma}

\begin{lemma}[Segment Division and Plane Separation]
    \label{lem:Plane_Separation_Consequence}
    Let A and B be distinct points in a plane, and let a transversal line $l$ intersect the line AB at a point F. Then F lies between A and B if and only if A and B lie on opposite sides of the line $l$. Consequently, if F is outside the segment AB, then A and B lie on the same side of $l$.
\end{lemma}

\begin{proof}
    This is a consequence of the Plane Separation Postulate, which states that any line in a plane divides the plane into two disjoint convex sets (half-planes). If two points lie in different half-planes, the segment connecting them must intersect the line. If F is between A and B, the segment AB contains F, which is on line $l$. Since F is the only intersection point of line AB and line $l$, A and B cannot both be in the same half-plane, as that would imply the segment AB is entirely contained within that half-plane, a contradiction. Thus, A and B are on opposite sides. Conversely, if A and B are on opposite sides of $l$, the segment AB must intersect $l$ at some point, which must be F. Therefore, F is on the segment AB.
\end{proof}

\begin{lemma}[Segment Ratio by Signed Perpendiculars]
    \label{lem:Segment_Ratio_by_Signed_Perpendiculars}
    \uses{def:Signed_Length}
    \uses{lem:AA_Similarity}
    \uses{lem:Plane_Separation_Consequence}
    Let A and B be two distinct points, and let a line $l$ intersect the line AB at a point F. Let $h_A$ and $h_B$ be the signed distances from A and B to the line $l$, where distances on opposite sides of $l$ have opposite signs. Then, the ratio of the signed lengths is given by:
    $$ \frac{\overline{AF}}{\overline{FB}} = -\frac{h_A}{h_B} $$
\end{lemma}

\begin{proof}
    Let $P_A$ and $P_B$ be the feet of the perpendiculars from A and B to line $l$. The triangles $\triangle AFP_A$ and $\triangle BFP_B$ are right-angled triangles. The angles $\angle AFP_A$ and $\angle BFP_B$ are vertically opposite, so they are equal. By lemma \ref{lem:AA_Similarity}, $\triangle AFP_A \sim \triangle BFP_B$. Therefore, the ratio of the absolute lengths of corresponding sides is equal:
    $$ \frac{|\overline{AF}|}{|\overline{FB}|} = \frac{|h_A|}{|h_B|} $$
    Now we analyze the signs.
    \begin{itemize}
        \item Case 1: F is between A and B. By definition \ref{def:Signed_Length}, the ratio $\frac{\overline{AF}}{\overline{FB}}$ is positive. By lemma \ref{lem:Plane_Separation_Consequence}, A and B are on opposite sides of $l$, so their signed distances $h_A$ and $h_B$ have opposite signs, which means the ratio $\frac{h_A}{h_B}$ is negative.
        \item Case 2: F is outside the segment AB. By definition \ref{def:Signed_Length}, the ratio $\frac{\overline{AF}}{\overline{FB}}$ is negative. By lemma \ref{lem:Plane_Separation_Consequence}, A and B are on the same side of $l$, so their signed distances $h_A$ and $h_B$ have the same sign, which means the ratio $\frac{h_A}{h_B}$ is positive.
    \end{itemize}
    In both cases, the ratio of signed lengths $\frac{\overline{AF}}{\overline{FB}}$ has the opposite sign to the ratio of signed distances $\frac{h_A}{h_B}$. Since their magnitudes are equal, we conclude:
    $$ \frac{\overline{AF}}{\overline{FB}} = -\frac{h_A}{h_B} $$
\end{proof}

\begin{theorem}[Menelaus's Theorem]
    \label{thm:Menelaus}
    \uses{def:Signed_Length}
    \uses{def:Transversal}
    \uses{lem:Segment_Ratio_by_Signed_Perpendiculars}
    Let $\triangle ABC$ be a triangle, and let a line be a transversal that intersects the lines containing the sides AB, BC, and CA at points F, D, and E, respectively. Then,
    $$ \frac{\overline{AF}}{\overline{FB}} \cdot \frac{\overline{BD}}{\overline{DC}} \cdot \frac{\overline{CE}}{\overline{EA}} = -1 $$
\end{theorem}

\begin{proof}
    \uses{lem:Segment_Ratio_by_Signed_Perpendiculars}
    Let the transversal line be denoted by $l$. Let $h_A$, $h_B$, and $h_C$ be the signed distances from the vertices A, B, and C to the line $l$.
    By applying lemma \ref{lem:Segment_Ratio_by_Signed_Perpendiculars} to each side of the triangle, we obtain the following three relations:
    \begin{align*}
        \frac{\overline{AF}}{\overline{FB}} &= -\frac{h_A}{h_B} \\
        \frac{\overline{BD}}{\overline{DC}} &= -\frac{h_B}{h_C} \\
        \frac{\overline{CE}}{\overline{EA}} &= -\frac{h_C}{h_A}
    \end{align*}
    Multiplying these three equations together yields:
    $$ \frac{\overline{AF}}{\overline{FB}} \cdot \frac{\overline{BD}}{\overline{DC}} \cdot \frac{\overline{CE}}{\overline{EA}} = \left(-\frac{h_A}{h_B}\right) \cdot \left(-\frac{h_B}{h_C}\right) \cdot \left(-\frac{h_C}{h_A}\right) = - \frac{h_A h_B h_C}{h_B h_C h_A} = -1 $$
    This completes the proof.
\end{proof}

\begin{lemma}[Uniqueness of a Point Dividing a Segment in a Given Ratio]
    \label{lem:Uniqueness_of_Division_Point}
    \uses{def:Signed_Length}
    Let A and B be two distinct points on a line. For any real number $k \neq -1$, there exists a unique point F on the line AB such that $\frac{\overline{AF}}{\overline{FB}} = k$.
\end{lemma}

\begin{proof}
    Let the line be represented by the real number line, with the coordinate of A being $a$ and the coordinate of B being $b$. Let the coordinate of F be $x$. The signed lengths are $\overline{AF} = x-a$ and $\overline{FB} = b-x$. The condition is:
    $$ \frac{x-a}{b-x} = k $$
    We solve for $x$:
    $$ x - a = k(b - x) \implies x - a = kb - kx \implies x(1+k) = a + kb $$
    Since $k \neq -1$, the term $(1+k)$ is non-zero, and we can divide by it to find a unique value for $x$:
    $$ x = \frac{a+kb}{1+k} $$
    This shows that for any $k \neq -1$, there is a unique point F that satisfies the condition.
\end{proof}

\begin{lemma}[Non-parallelism Condition for Converse Menelaus]
    \label{lem:Non_Parallelism}
    \uses{def:Signed_Length}
    Let $\triangle ABC$ be a triangle and let D, E, F be points on lines BC, CA, AB respectively, satisfying the Menelaus condition $\frac{\overline{AF}}{\overline{FB}} \cdot \frac{\overline{BD}}{\overline{DC}} \cdot \frac{\overline{CE}}{\overline{EA}} = -1$. Then the line DE cannot be parallel to the line AB.
\end{lemma}

\begin{proof}
    \uses{def:Signed_Length}
    Assume for contradiction that the line DE is parallel to the line AB. By the property of similar triangles ($\triangle CDE \sim \triangle CBA$), we have $\frac{\overline{CE}}{\overline{EA}} = \frac{\overline{CD}}{\overline{DB}}$.
    This can be rewritten as $\frac{\overline{CE}}{\overline{EA}} = \frac{-\overline{DC}}{-\overline{BD}} = \frac{\overline{DC}}{\overline{BD}}$.
    Therefore, $\frac{\overline{BD}}{\overline{DC}} \cdot \frac{\overline{CE}}{\overline{EA}} = 1$.
    Substituting this into the given Menelaus condition gives:
    $$ \frac{\overline{AF}}{\overline{FB}} \cdot \left( \frac{\overline{BD}}{\overline{DC}} \cdot \frac{\overline{CE}}{\overline{EA}} \right) = -1 \implies \frac{\overline{AF}}{\overline{FB}} \cdot (1) = -1 $$
    This implies $\frac{\overline{AF}}{\overline{FB}} = -1$, which means $\overline{AF} = -\overline{FB}$. By the property of signed lengths, $\overline{AF} = \overline{AB} + \overline{BF} = \overline{AB} - \overline{FB}$.
    Substituting $\overline{AF} = -\overline{FB}$ gives $\overline{AB} - \overline{FB} = -\overline{FB}$, which simplifies to $\overline{AB} = 0$. This implies that A and B are the same point, which contradicts the definition of a triangle.
    Thus, our initial assumption that DE is parallel to AB must be false.
\end{proof}

\begin{theorem}[Converse of Menelaus's Theorem]
    \label{thm:Converse_of_Menelaus}
    \uses{def:Signed_Length}
    \uses{thm:Menelaus}
    \uses{lem:Uniqueness_of_Division_Point}
    \uses{lem:Non_Parallelism}
    Let $\triangle ABC$ be a triangle, and let D, E, and F be points on the lines containing the sides BC, CA, and AB, respectively. If
    $$ \frac{\overline{AF}}{\overline{FB}} \cdot \frac{\overline{BD}}{\overline{DC}} \cdot \frac{\overline{CE}}{\overline{EA}} = -1 $$
    then the points D, E, and F are collinear.
\end{theorem}

\begin{proof}
    By lemma \ref{lem:Non_Parallelism}, the line DE is not parallel to the line AB. Therefore, the line passing through D and E must intersect the line AB at exactly one point. Let's call this point F'.
    Since D, E, and F' are collinear, by theorem \ref{thm:Menelaus}, they must satisfy the equation:
    $$ \frac{\overline{AF'}}{\overline{F'B}} \cdot \frac{\overline{BD}}{\overline{DC}} \cdot \frac{\overline{CE}}{\overline{EA}} = -1 $$
    We are given that the point F satisfies:
    $$ \frac{\overline{AF}}{\overline{FB}} \cdot \frac{\overline{BD}}{\overline{DC}} \cdot \frac{\overline{CE}}{\overline{EA}} = -1 $$
    Since D and E are not vertices, the terms $\frac{\overline{BD}}{\overline{DC}}$ and $\frac{\overline{CE}}{\overline{EA}}$ are well-defined, finite, and non-zero. Comparing the two equations gives:
    $$ \frac{\overline{AF'}}{\overline{F'B}} = \frac{\overline{AF}}{\overline{FB}} $$
    Let this ratio be $k$. Lemma \ref{lem:Non_Parallelism} ensures that $k \neq -1$.
    By lemma \ref{lem:Uniqueness_of_Division_Point}, there is only one point on the line AB that divides the segment AB in a given ratio $k$. Therefore, F and F' must be the same point.
    Since F' is on the line DE by construction, F must also be on the line DE. Thus, the points D, E, and F are collinear.
\end{proof} 